\chapter{Moral Hazard and Contracts}\label{ch:b2:moral-hazard}

In 1994 and 1995 Safelite, the largest installer of car windscreens in the United States, moved its installers from
hourly wages to pay by the piece, region by region, with a guarantee of about the former wage. The firm kept monthly
records of each installer's output before and after. The chart shows three estimates of how much output per worker
rose.

\begin{openingchart}{Rise in windscreens installed per worker per day after Safelite moved from hourly wages to piece
  rates, per cent: raw means of all monthly observations (2.70 units a day under hourly wages, 3.24 under piece rates);
  with month and year effects controlled; and within the same worker, comparing each installer with himself.}{Lazear
  (2000), Tables~2--3 and text}
\begin{tikzpicture}
\begin{axis}[omrecord, xbar, bar width=10pt, width=0.62\linewidth, height=0.36\linewidth, xmin=0, xmax=52,
    ymin=0.4, ymax=3.6, ymajorgrids=false, xlabel={Per cent},
    ytick=data, yticklabels from table={figdata/bachelor-2/17-moral-hazard/out/safelite.csv}{label},
    nodes near coords, nodes near coords style={font=\footnotesize}]
  \addplot[fill=omThm!50, draw=omThm] table[x=pct, y=pos]
    {figdata/bachelor-2/17-moral-hazard/out/safelite.csv};
\end{axis}
\end{tikzpicture}
\end{openingchart}

% ledger: b2:lazear2000-design, b2:lazear2000-t2, b2:lazear2000-t3
% figdata: safelite_summary (raw 20, total 44, within 22, between 22, within_share 0.5)
Output per worker rose by 44 per cent once month and year effects are controlled. Half of that, 22 per cent, is the same
installer working faster: an incentive effect. The other half came from who worked there, as the firm kept more of its
most productive installers and hired more productive ones. Pay by results changed both what each worker did and who
chose to work.

Under an hourly wage, the firm could not see each installer's effort well enough to pay for it. That is the problem of
this chapter: one party acts for another, and the other cannot observe the action. The chapter builds the model of
contracts written under that constraint, finds what the best contract trades off, and turns to the evidence from
workshops, farms and health insurance. It ends with contracts that cannot be written at all.

\begin{recall}
From \cref{ch:b2:expected-utility,ch:b2:insurance-and-risk-sharing}: expected utility, the \omterm{def:b2:expected-utility:risk-premium}{risk premium}, and insurance
against risk. From \cref{ch:b2:adverse-selection,ch:b2:signalling-and-screening}: hidden information, participation and
\omterm{def:b2:signalling-and-screening:ic}{incentive-compatibility constraints}.
\end{recall}

\begin{notation}[Symbols of this chapter]\label{not:b2:moral-hazard:symbols}
The agent's effort is $e$ and its cost $\mathit{C}(e)$; output is $y$ and the agent's pay $w$. In the linear model, $w = a + by$, with $b$ the
share of output paid to the agent, $r$ her \omterm{def:b2:expected-utility:ara}{coefficient of absolute risk aversion}, $\sigma^2$ the variance of output given effort,
and $c$ the slope of her marginal cost of effort. The agent's reservation utility is $\bar u$.
\end{notation}

\section{Principal and agent}\label{sec:b2:moral-hazard:principal-agent}

\begin{definition}[Principal--agent problem]\label{def:b2:moral-hazard:principal-agent}
A \emph{principal--agent problem}\index{principal--agent problem} arises when one party, the principal, has an
action taken on its behalf by another, the agent, whose interests differ from its own, and must design the agent's
reward without full information about the agent's action or knowledge.
\end{definition}

\begin{definition}[Moral hazard]\label{def:b2:moral-hazard:moral-hazard}
\emph{Moral hazard}\index{moral hazard} is the problem of a hidden action: after a contract is signed, one party takes
an action that affects the other's payoff and that the other cannot observe or cannot write into the contract.
\end{definition}

\omterm{def:b2:moral-hazard:moral-hazard}{Moral hazard} differs from \omterm{def:b2:adverse-selection:adverse-selection}{adverse selection} (\cref{ch:b2:adverse-selection}) in timing: the information is hidden before
the contract in \omterm{def:b2:adverse-selection:adverse-selection}{adverse selection}, and the action is hidden after it here. An insurer may face both: people who know they
are risky buy more cover, and people with cover take less care.

Take an orchard owner, the principal, who hires a manager. The manager can work hard, at a cost to her of 2 (in the units
of her utility), or not, at no cost. Output, in thousands of dalens, is 100 or 20; hard work makes 100 likely with
probability 0.8, easy work with probability 0.4. The manager's utility is $\sqrt{w} - \text{cost}$, and she will take the job only
if it gives her at least $\bar u = 5$, what she could earn elsewhere. The owner is risk neutral.

If the owner could see effort, he would pay a fixed wage for hard work, the smallest that the manager accepts: $\sqrt{w} - 2
= 5$, so $w = 49$. Expected output is $0.8 \times 100 + 0.2 \times 20 = 84$, and profit 35. Easy work would need only $w = 25$, with expected output
52 and profit 27. So the first best is hard work, at a fixed wage: the risk-neutral owner bears all the risk.

\section{Insurance against incentives}\label{sec:b2:moral-hazard:second-best}

\begin{definition}[Second-best contract]\label{def:b2:moral-hazard:second-best}
A \emph{second-best contract}\index{second-best contract} is the contract that maximises the principal's payoff
subject to the agent's participation and incentive constraints when the agent's action cannot be contracted on.
\end{definition}

If effort is hidden, a fixed wage gives the manager no reason to work hard. The owner must pay more when output is high,
$w_H$, than when it is low, $w_L$. Writing $u_H = \sqrt{w_H}$ and $u_L = \sqrt{w_L}$, hard work must be worth its cost (incentive
compatibility) and the job must be worth taking (participation):
\[
  0.8u_H + 0.2u_L - 2 \ge 0.4u_H + 0.6u_L, \qquad 0.8u_H + 0.2u_L - 2 \ge 5 .
\]
The first says $u_H - u_L \ge 5$. The owner wants the cheapest such pair, and both bind: $u_H - u_L = 5$ and $0.8u_H + 0.2u_L = 7$, so $u_L = 3$ and
$u_H = 8$, that is $w_L = 9$ and $w_H = 64$. The expected wage is $0.8 \times 64 + 0.2 \times 9 = 53$, against 49 in the first best.

\begin{omfigure}
\begin{tikzpicture}[econ]
\begin{axis}[omrecord, grid=none, xmin=0, xmax=10, ymin=0, ymax=12, width=0.5\linewidth,
    height=0.5\linewidth, xlabel={$u_L$: utility of pay if output is low}, ylabel={$u_H$: if high},
    xtick={3,7}, ytick={7,8}, legend style={at={(1.05,0.5)}, anchor=west}]
  \addplot[black!50, thin, domain=0:10, samples=2] {x};
  \addplot[omDef, thick, domain=0:8.75, samples=2] {(7 - 0.2*x)/0.8};
  \addplot[omThm, thick, dashed, domain=0:7, samples=2] {x + 5};
  \node[eqpoint] at (axis cs:3,8) {};
  \node[eqpoint] at (axis cs:7,7) {};
  \node[anchor=south east, font=\footnotesize] at (axis cs:2.9,8.9) {second best};
  \node[anchor=north west, font=\footnotesize] at (axis cs:7.2,6.5) {first best};
  \legend{Equal pay, Participation, Incentive}
\end{axis}
\end{tikzpicture}

{\small The two-outcome contract in utilities. The first best pays the same in both outcomes, on the equal-pay line.
With hidden effort the contract must lie on or above the incentive line, five units of utility apart; the cheapest point
that also meets participation is $u_L = 3$, $u_H = 8$.}
\end{omfigure}

The manager is no worse off, at expected utility 7 minus the cost of 2. The 4 thousand extra is the price of risk: to make
her work, the owner must expose her to risk, and a risk-averse manager must be paid to bear it. That is the
\emph{agency cost} of hidden effort. Profit falls from 35 to 31, still above the 27 of easy work, so the owner still asks
for hard work. Were the manager risk neutral, the cost would vanish: making her bear all the risk costs nothing to insure
(\cref{exo:b2:moral-hazard:11}). The trade-off is between insurance and incentives.

\section{Linear contracts}\label{sec:b2:moral-hazard:linear}

\begin{definition}[Linear incentive contract]\label{def:b2:moral-hazard:linear-contract}
A \emph{linear incentive contract}\index{linear incentive contract} pays the agent a fixed amount plus a constant share
of a measure of output: $w = a + by$.
\end{definition}

\begin{definition}[Piece rate]\label{def:b2:moral-hazard:piece-rate}
A \emph{piece rate}\index{piece rate} is pay by the unit of output produced, a linear contract on a count of units.
\end{definition}

Contracts in practice are often linear: commissions, \omterm{def:b2:moral-hazard:piece-rate}{piece rates}, crop shares. A model in which the linear contract is
the right one to study takes output $y = e + \varepsilon$, with $\varepsilon$ normal with mean zero and variance $\sigma^2$; effort cost $ce^2/2$; and an agent
with \omterm{def:b2:expected-utility:cara}{constant absolute risk aversion} $r$. For such an agent a normal income with mean $\mu$ and variance $v$ is worth its
\omterm{def:b2:expected-utility:certainty-equivalent}{certainty equivalent} $\mu - rv/2$. Under $w = a + by$ her \omterm{def:b2:expected-utility:certainty-equivalent}{certainty equivalent} is
\[
  a + be - \frac{c e^2}{2} - \frac{r b^2 \sigma^2}{2},
\]
so she chooses $e = b/c$: effort rises with the share. The fixed part $a$ does not affect effort; it only divides the surplus,
and the principal sets it so that participation just binds.

\begin{proposition}[The optimal linear share]\label{prop:b2:moral-hazard:linear}
In the linear model the share that maximises the principal's expected profit is
\[
  b^* = \frac{1}{1 + rc\sigma^2},
\]
and the total certainty-equivalent surplus is $b^*/(2c)$, against $1/(2c)$ when effort is observable.
\end{proposition}

\begin{proof}
With participation binding, the principal's profit is the total surplus, expected output less the cost of effort and the
agent's \omterm{def:b2:expected-utility:risk-premium}{risk premium}, minus the fixed reservation level. With $e = b/c$, the surplus is
\[
  S(b) = \frac{b}{c} - \frac{b^2}{2c} - \frac{r b^2\sigma^2}{2}.
\]
Its derivative $1/c - b/c - rb\sigma^2$ is zero at $b^* = 1/(1 + rc\sigma^2)$, and $S$ is concave, so $b^*$ is the maximum. Since $1/c + r\sigma^2 = 1/(cb^*)$, $S(b^*) =
b^*/c - b^{*2}/(2cb^*) = b^*/(2c)$. With observable effort, the principal asks for $e = 1/c$ and pays a fixed wage, so the surplus is
$1/c - 1/(2c) = 1/(2c)$.
\end{proof}

The share falls with \omterm{def:b2:expected-utility:risk-aversion}{risk aversion}, with the cost of effort and with noise: when output says little about effort, paying
on output buys little effort at a high price in risk. The listing finds the best share by search and checks the formula.

\omcode{code/bachelor-2/17-moral-hazard/listings/linear.py}{1}{25}{The best linear share for six levels of noise, found by
search and by the formula.}\label{lst:b2:moral-hazard:linear}

\omoutput{code/bachelor-2/17-moral-hazard/outputs/linear.txt}{The share, effort and surplus fall as output gets noisier;
the agency cost rises toward the whole first-best surplus, 0.5.}\label{out:b2:moral-hazard:linear}

% code: listings/linear.py
\begin{omfigure}
\begin{tikzpicture}[econ]
\begin{axis}[omrecord, grid=none, xmin=0, xmax=4.2, ymin=0, ymax=1.05, width=0.6\linewidth,
    height=0.4\linewidth, xlabel={$\sigma^2$: variance of output}, ylabel={Optimal share $b^*$},
    xtick={0,1,2,3,4}, ytick={0,0.5,1}, legend style={at={(1.03,0.5)}, anchor=west}]
  \addplot[omDef, thick, domain=0:4.2, samples=60] {1/(1 + 1*x)};
  \addplot[omThm, thick, dashed, domain=0:4.2, samples=60] {1/(1 + 2*x)};
  \addplot[omMeth, thick, dotted, domain=0:4.2, samples=60] {1/(1 + 4*x)};
  \legend{$rc = 1$, $rc = 2$, $rc = 4$}
\end{axis}
\end{tikzpicture}

{\small The risk--incentive trade-off. The optimal share of output paid to the agent falls as output gets noisier, and
faster for a more risk-averse agent or one whose effort is more costly.}
\end{omfigure}

\section{Several tasks}\label{sec:b2:moral-hazard:multitask}

Most jobs have several tasks, and some are easier to measure than others. A teacher raises test scores and also teaches
what tests do not measure; an installer fits windscreens quickly and also fits them well. Suppose the agent's cost
depends on her total effort, and only the first task is measured. Pay on the measure moves effort from the second task to
the first, at no extra cost to her. If the principal values both, strong pay on the measured task can lower what he gets.
The model's answer is weaker incentives than one task alone would warrant, or a fixed wage with supervision, when what is
measured is a poor guide to what is wanted. Safelite could pay by the piece in part because a badly fitted windscreen soon came
back, and the installer who fitted it had to refit it on his own time (Lazear 2000).

\section{Teams and relative performance}\label{sec:b2:moral-hazard:teams}

When output is joint, pay on it is shared. If $n$ workers each choose effort $e_i$ at cost $e_i^2/2$, output is $\sum_i e_i$ and is
shared equally, each worker earns $1/n$ of what her own effort adds, and chooses $e_i = 1/n$; the efficient effort is 1. A team of
four works at a quarter of the efficient level. This is free riding, as with the public goods of the school volumes.

\begin{definition}[Relative performance evaluation]\label{def:b2:moral-hazard:rpe}
\emph{Relative performance evaluation}\index{relative performance evaluation} is pay that depends on an agent's
output compared with the output of others who face some of the same shocks.
\end{definition}

Comparison helps when agents share a shock. Let $y_1 = e_1 + \eta + \varepsilon_1$ and $y_2 = e_2 + \eta + \varepsilon_2$, where $\eta$ is a common shock with variance
$\sigma_\eta^2$, such as the weather on two neighbouring farms or the state of the market for two salespeople, and $\varepsilon_1$, $\varepsilon_2$ are
independent with variance $\sigma^2$. Pay agent 1 on $y_1 - k y_2$. Since $e_2$ is expected, subtracting $ky_2$ moves the measure by a known amount,
which the fixed part can offset; it changes only the noise.

\begin{proposition}[Relative evaluation filters a common shock]\label{prop:b2:moral-hazard:rpe}
The variance of $y_1 - ky_2$ is least at $k^* = \sigma_\eta^2/(\sigma_\eta^2 + \sigma^2)$, where it equals $\sigma^2 + \sigma_\eta^2\sigma^2/(\sigma_\eta^2 + \sigma^2)$, below the variance $\sigma^2 + \sigma_\eta^2$ of $y_1$
whenever $\sigma_\eta^2 > 0$. So the optimal linear contract pays on $y_1 - k^*y_2$, with a larger share $b^*$.
\end{proposition}

\begin{proof}
$y_1 - ky_2 - (e_1 - ke_2) = (1 - k)\eta + \varepsilon_1 - k\varepsilon_2$, with variance $V(k) = (1 - k)^2\sigma_\eta^2 + \sigma^2 + k^2\sigma^2$. Then $V'(k) = -2(1 - k)\sigma_\eta^2 +
2k\sigma^2$, zero at $k^*$, and $V$ is convex. Substituting, $V(k^*) = \sigma^2 + \sigma_\eta^2\sigma^2/(\sigma_\eta^2 + \sigma^2)$. Since the second term is below $\sigma_\eta^2$, $V(k^*) < V(0)$. By
\cref{prop:b2:moral-hazard:linear}, a lower variance gives a larger share and a larger surplus.
\end{proof}

\begin{omfigure}
\begin{tikzpicture}[econ]
\begin{axis}[omrecord, grid=none, xmin=0, xmax=1.25, ymin=0, ymax=4.5, width=0.55\linewidth,
    height=0.4\linewidth, xlabel={$k$: weight on the other agent's output}, ylabel={Variance of the measure},
    xtick={0,0.75,1}, ytick={1,1.75,2,4}]
  \addplot[omDef, thick, domain=0:1.25, samples=60] {3*(1 - x)^2 + 1 + x^2};
  \addplot[black!50, dashed, domain=0:1.25, samples=2] {1};
  \node[eqpoint] at (axis cs:0.75,1.75) {};
  \node[anchor=north, font=\footnotesize] at (axis cs:0.75,1.7) {$k^*$};
  \node[anchor=south west, font=\footnotesize] at (axis cs:0,4) {own output};
  \node[anchor=south, font=\footnotesize] at (axis cs:0.55,0.95) {own idiosyncratic noise};
\end{axis}
\end{tikzpicture}

{\small Filtering a common shock, with $\sigma_\eta^2 = 3$ and $\sigma^2 = 1$. Paying on own output alone carries a variance of 4; on the difference
from the other agent's output, 2; on $y_1 - 0.75y_2$, 1.75. No weight brings it to the idiosyncratic noise alone, since the other
agent's output carries noise of its own.}
\end{omfigure}

Relative evaluation has costs: agents may collude to work less, sabotage one another, or choose colleagues badly; and
if the shock is not common it only adds noise.

\section{The evidence}\label{sec:b2:moral-hazard:evidence}

\begin{record}[Piece rates at Safelite]
% ledger: b2:lazear2000-t2, b2:lazear2000-t3
Output per installer per day was 2.70 units under hourly wages and 3.24 under \omterm{def:b2:moral-hazard:piece-rate}{piece rates}; its spread across installers
widened, with a standard deviation of 1.42 before and 1.59 after. Pay rose little, from 2{,}228 to 2{,}283 dollars a month,
while the cost of an installed unit fell from 44.43 to 35.24 dollars. The model's predictions hold here: output rose, more
productive workers stayed or came, and the gain was shared between the firm and its workers (Lazear 2000).
\end{record}

\begin{model}{Principal and agent with linear contracts}
\assumes A risk-averse agent whose effort raises output; output observed with noise; effort not observable or not
contractible; the principal chooses a fixed part and a share.
\predicts Effort and output rise with the share paid on output; the share falls with noise, \omterm{def:b2:expected-utility:risk-aversion}{risk aversion} and the cost
of effort; comparison with others who share the same shocks raises it.
\tested \omterm{def:b2:moral-hazard:piece-rate}{Piece rates} at Safelite (Lazear 2000); higher output after a reform that raised sharecroppers' shares in West Bengal
(Banerjee, Gertler and Ghatak 2002); medical use falling with the share of costs the insured pay (Brook et al.\ 2006).
\breaks When the measure can be gamed at the expense of what is not measured, the several-task problem above; and when pay
changes why people work as well as how hard, which this model leaves out.
\end{model}

\begin{record}[Tenancy reform in West Bengal]
% ledger: b2:bgg2002
A sharecropper pays the landowner a share of the crop: a linear contract with share $b$ equal to what the tenant keeps. Within a year of its election in 1977, the West Bengal state government began to register sharecroppers, giving the registered permanent, inheritable
tenure as long as they paid the landowner the legal share, at most 25 per cent of output. By 1993 about 65 per cent of
sharecroppers were registered, against 15 per cent before. In a survey, the share of tenants keeping more than half of
their crop rose from 17 to 39 per cent. From 1979 to 1993 rice yields rose 69 per cent in West Bengal and 44 per cent in
neighbouring Bangladesh, after growing at similar rates over the decade before. With controls for rainfall, irrigation and
high-yield varieties, the authors estimate that yields were 18 per cent higher in West Bengal in 1988--91, and that the reform
explains about 28 per cent of the later growth of farm productivity there (Banerjee, Gertler and Ghatak 2002). The reform changed
both the share and the security of tenure; the study cannot separate them.
\end{record}

\begin{definition}[Coinsurance rate]\label{def:b2:moral-hazard:coinsurance}
The \emph{coinsurance rate}\index{coinsurance rate} of an insurance policy is the share of each covered cost that the
insured person pays.
\end{definition}

Health insurance lowers the price of care at the point of use, so the insured may use more of it. That is \omterm{def:b2:moral-hazard:moral-hazard}{moral hazard}
too, though the hidden action is a choice of how much care to buy. To measure it one must compare people who differ in
cover but not in health, which is what an experiment can do.

\begin{omchart}{Mean outpatient visits to a doctor per person-year by the coinsurance rate of the randomly assigned
  plan, RAND Health Insurance Experiment, plans without an individual deductible (14{,}941 person-years). Whiskers: two
  standard errors.}{statsmodels bundled data (public domain), randhie, as distributed by Cameron and Trivedi; our
  computation}
\begin{tikzpicture}
\begin{axis}[omrecord, ybar, bar width=16pt, width=0.6\linewidth, height=0.42\linewidth, xmin=-12, xmax=107,
    ymin=0, ymax=4.2, xtick={0,25,50,95}, xlabel={Coinsurance rate (per cent)},
    ylabel={Doctor visits a year}, xmajorgrids=false]
  \addplot[fill=omThm!50, draw=omThm, error bars/.cd, y dir=both, y explicit]
    table[x=coins, y=mean, y error expr=2*\thisrow{se}]
    {figdata/bachelor-2/17-moral-hazard/out/randhie.csv};
\end{axis}
\end{tikzpicture}
\end{omchart}

\begin{record}[The RAND Health Insurance Experiment]
% ledger: b2:randhie2006
% figdata: randhie_summary (free 3.555, c25 2.787, c50 2.561, c95 2.112, fall_pct 40.6, n 14941)
From 1971 to 1982, 2{,}750 families, more than 7{,}700 people under 65 at six sites, were assigned at random to free care or
to plans with coinsurance of 25, 50 or 95 per cent, with out-of-pocket spending capped. Those with 25 per cent coinsurance
spent 20 per cent less than those with free care, by using fewer services, not by paying lower prices. In general the
reduction had no adverse effect on health, but free care improved blood pressure, dental health, vision and some serious
symptoms, mainly among the sickest and poorest (Brook et al.\ 2006). In the experiment's public data
(the chart above), visits fall from 3.56 a year under free care to 2.11 under 95 per cent coinsurance, 41 per cent
fewer.
\end{record}

\section{Contracts that cannot be written}\label{sec:b2:moral-hazard:incomplete}

\begin{definition}[Incomplete contract]\label{def:b2:moral-hazard:incomplete-contract}
An \emph{incomplete contract}\index{incomplete contract} is one that does not specify what each party must do in every
possible situation, because some situations cannot be foreseen, described or verified by a court.
\end{definition}

\begin{definition}[Hold-up problem]\label{def:b2:moral-hazard:hold-up}
The \emph{hold-up problem}\index{hold-up problem} arises when one party makes an investment whose value depends on
trading with a particular partner, and the terms of trade are bargained after the investment: the partner can capture
part of its return, so the investment is too small or not made.
\end{definition}

A supplier can spend 30 thousand dalens adapting its plant to one buyer's design, which raises the value of their trade by
50. If the contract could specify the price after the investment, the supplier would invest. If it cannot, because the
design and quality cannot be described in advance, the price is bargained after the plant is adapted, and the investment
cost is then sunk. If they split the gain equally, the supplier gets 25 for its 30 and does not invest, though the
investment would create 20 net. Remedies include long-term contracts, joint ownership, and one firm owning the other: which
activities a firm does itself and which it buys can be explained in part by which investments would otherwise be held up.
Each remedy has costs of its own, such as weaker incentives inside a larger firm.

\section{Exercises}\label{sec:b2:moral-hazard:exercises}

\begin{exercise}[$\star$]\label{exo:b2:moral-hazard:1}
Read the opening chart. Why is the raw gain smaller than the gain with month and year controlled, and what does the gap
between the second and third bars measure?
\end{exercise}

\begin{exercise}[$\star$]\label{exo:b2:moral-hazard:2}
In the orchard of \cref{sec:b2:moral-hazard:principal-agent}, let the manager's outside option be $\bar u = 6$. Find the first-best wage and
profit, and the profit from easy work.
\end{exercise}

\begin{exercise}[$\star$]\label{exo:b2:moral-hazard:3}
From the chart of the RAND experiment in \cref{sec:b2:moral-hazard:evidence}, compute the fall in visits, in per cent, from free care to 25 per cent coinsurance, and from 25
to 95 per cent. Why does random assignment matter for reading these as effects of the plan?
\end{exercise}

\begin{exercise}[$\star$]\label{exo:b2:moral-hazard:4}
Predict, then compute: in a team of four sharing output equally, with cost $e^2/2$, what effort does each choose, and what is
each one's surplus, against the efficient level?
\end{exercise}

\begin{exercise}[$\star\star$]\label{exo:b2:moral-hazard:5}
With $\bar u = 6$, find the \omterm{def:b2:moral-hazard:second-best}{second-best contract}, the expected wage and the agency cost. Does the owner still want hard work?
\end{exercise}

\begin{exercise}[$\star\star$]\label{exo:b2:moral-hazard:6}
In the linear model with $r = 2$, $c = 1$ and $\sigma^2 = 0.5$, find $b^*$, effort, the agent's \omterm{def:b2:expected-utility:risk-premium}{risk premium}, the cost of effort, the surplus and
the agency cost.
\end{exercise}

\begin{exercise}[$\star\star$]\label{exo:b2:moral-hazard:7}
In the hold-up example, find the smallest share of the gain that makes the supplier invest. Why might a contract fixing
the price in advance not solve the problem?
\end{exercise}

\begin{exercise}[$\star\star$]\label{exo:b2:moral-hazard:8}
Explain why the RAND results measure \omterm{def:b2:moral-hazard:moral-hazard}{moral hazard} and not \omterm{def:b2:adverse-selection:adverse-selection}{adverse selection}, and why the same comparison across ordinary
insurance buyers would not.
\end{exercise}

\begin{exercise}[$\star\star$]\label{exo:b2:moral-hazard:9}
With $\sigma_\eta^2 = 3$ and $\sigma^2 = 1$, compute the variance of the measure for $k = 0$, $k = 1$ and $k^*$. With $r = 2$ and $c = 1$, find the optimal share in
each case.
\end{exercise}

\begin{exercise}[$\star\star\star$]\label{exo:b2:moral-hazard:10}
Show that the agency cost in the linear model is $(1 - b^*)/(2c)$, and that it rises with $\sigma^2$ toward the whole first-best surplus.
\end{exercise}

\begin{exercise}[$\star\star\star$]\label{exo:b2:moral-hazard:11}
In the two-outcome model, let the manager be risk neutral, with utility $w - \text{cost}$. Find the cheapest contract for hard work
and show that the agency cost is zero. Interpret.
\end{exercise}

\begin{exercise}[$\star\star\star$]\label{exo:b2:moral-hazard:12}
Agent 1 is compared with the average output of $n$ others, each with idiosyncratic variance $\sigma^2$, all sharing the shock $\eta$. Find
the best weight $k$ on the average and the variance of the measure. Evaluate for $\sigma_\eta^2 = 3$, $\sigma^2 = 1$ and $n = 4$, and find the limit as $n$ grows.
\end{exercise}

\section{Problem: Sharecropping in a Brenland orchard}\label{sec:b2:moral-hazard:problem}

\begin{problem}[{Weekend problem --- an orchard owner in Brenland lets his orchard to a tenant for a share of the crop: what
share is best, what does hidden effort cost, and how much does comparing with the neighbouring orchard save}]
\label{pb:b2:moral-hazard:1}
The crop is worth $y = e + \eta + \varepsilon$ hundred dalens a season, where $e$ is the tenant's effort, $\eta$ the weather, common to the
neighbouring orchard, with variance 1{,}800, and $\varepsilon$ the orchard's own luck, with variance 200. Effort costs the tenant $0.01e^2/2$. The
tenant keeps a share $b$ of the crop plus a fixed payment $a$ (negative if he pays a rent), and has \omterm{def:b2:expected-utility:cara}{constant absolute risk
aversion} $r = 0.05$. The owner is risk neutral. Effort cannot be observed.

\medskip
\noindent\textbf{Part I --- Effort and output.}
\begin{enumerate}
  \item What effort does the tenant choose for a share $b$?
  \item For $b = 0.5$, find expected output and the tenant's cost of effort.
  \item For $b = 0.5$, find the variance of the crop and the tenant's \omterm{def:b2:expected-utility:risk-premium}{risk premium}.
  \item Write the tenant's \omterm{def:b2:expected-utility:certainty-equivalent}{certainty equivalent} as a function of $a$ and $b$.
  \item Why does $a$ not affect effort?
\end{enumerate}

\medskip
\noindent\textbf{Part II --- The first best.}
\begin{enumerate}[resume]
  \item If the owner could observe effort, what effort would he ask for?
  \item What would the total surplus be?
  \item With a fixed wage and no share, what effort and surplus result?
  \item With a fixed rent, the tenant keeping the whole crop, what effort, \omterm{def:b2:expected-utility:risk-premium}{risk premium} and surplus result?
  \item Why is neither extreme best?
\end{enumerate}

\medskip
\noindent\textbf{Part III --- The linear share.}
\begin{enumerate}[resume]
  \item Write the surplus as a function of $b$.
  \item Find the optimal share.
  \item Find effort and expected output.
  \item Find the \omterm{def:b2:expected-utility:risk-premium}{risk premium}, the cost of effort and the surplus.
  \item What is the agency cost?
  \item A law fixes the tenant's share at 0.75. Find effort and the surplus.
  \item Why can such a law lower the surplus in this model, though in West Bengal output rose after a similar reform?
\end{enumerate}

\medskip
\noindent\textbf{Part IV --- The neighbour's yield.}
\begin{enumerate}[resume]
  \item What share of the crop's variance does the weather account for?
  \item The contract pays on the orchard's crop minus the neighbour's, whose own luck also has variance 200. Find the variance
    of this measure and the optimal share.
  \item Find effort, the surplus and the agency cost.
  \item Find the best weight on the neighbour's crop and the variance of the measure.
  \item Find the optimal share, effort, surplus and agency cost.
  \item With the best weight, what surplus does the 0.75 law now give?
  \item Name two ways in which comparing with the neighbour could go wrong.
  \item State the result: the optimal share, the agency cost, and its fall with relative evaluation.
\end{enumerate}
\end{problem}
